\subsection{Characterization of the support of a measure}

By definition, if the support of a function \(f \in \mathcal{K}(\mathrm{X};\mathbf{C})\) does not intersect the support of a measure \(\mu\) , then \(\mu(f) = 0\) ; but the following more precise result is true:

PROPOSITION 8. --- Let \(\mu\) be a measure on a locally compact space \(X\) . For every function \(f \in \mathcal{K}(X; \mathbf{C})\) that is zero on \(\operatorname{Supp}(\mu)\) , \(\mu(f) = 0\) .

Set \(\mathrm{K} = \operatorname{Supp}(f)\) , \(\mathrm{S} = \operatorname{Supp}(\mu)\) . Given a number \(\varepsilon > 0\) , let \(\mathrm{V}\) be the set of \(x \in \mathrm{X}\) such that \(|f(x)| < \varepsilon\) ; \(\mathrm{V}\) is an open set containing \(\mathrm{S}\) by hypothesis, therefore \(\complement \mathrm{S}\) is a neighborhood of the compact set \(\complement \mathrm{V}\) . It follows that there exists a continuous mapping \(h\) of \(\mathrm{X}\) into {[}0,1{]}, equal to 1 on \(\complement \mathrm{V}\) and with support contained in \(\complement \mathrm{S}\) (§1, No.~2, Lemma 1). Since the support of \(fh\) does not intersect \(\mathrm{S}\) , \(\mu(fh) = 0\) . On the other hand, \(f = fh\) on \(\mathrm{K} \cap \complement \mathrm{V}\) , and \(|fh| \leqslant |f|\) on \(\mathrm{X}\) , therefore \(|f - fh| \leqslant 2\varepsilon\) on \(\mathrm{X}\) , by the choice of \(\mathrm{V}\) . Observe finally that there exists a number \(\mathrm{M}_{\mathrm{K}}\) such that \(|\mu(g)| \leqslant \mathrm{M}_{\mathrm{K}} \| g\|\) for every function \(g \in \mathcal{K}(\mathrm{X};\mathbf{C})\) whose support is contained in \(\mathrm{K}\) ; since the support of \(f - fh\) is contained in \(\mathrm{K}\) , \(|\mu(f - fh)| \leqslant 2\mathrm{M}_{\mathrm{K}}\varepsilon\) and consequently \(|\mu(f)| = |\mu(f - fh)| \leqslant 2\mathrm{M}_{\mathrm{K}}\varepsilon\) ; since \(\varepsilon\) is arbitrary, \(\mu(f) = 0\) .

COROLLARY 1. --- If two functions \(f, g\) in \(\mathcal{K}(\mathrm{X};\mathbf{C})\) are equal on \(\operatorname{Supp}(\mu)\) , then \(\mu(f) = \mu(g)\) .

COROLLARY 2. --- Let \(\mu\) be a positive measure on X; if \(f \in \mathcal{K}(X; \mathbf{C})\) is such that \(f(x) \geqslant 0\) on \(\operatorname{Supp}(\mu)\) , then \(\mu(f) \geqslant 0\) .

For, \(f = |f|\) on \(\operatorname{Supp}(\mu)\) , therefore \(\mu(f) = \mu(|f|) \geqslant 0\) by Corollary 1.

COROLLARY 3. --- Let \(\mu\) be a bounded measure on X; if \(f \in \mathcal{K}(X; \mathbf{C})\) is such that \(|f(x)| \leqslant a\) on \(\operatorname{Supp}(\mu)\) , then \(|\mu(f)| \leqslant a \| \mu \|\) .

For, \(\operatorname{Supp}(|\mu|) = \operatorname{Supp}(\mu)\) , and if \(h\) is a continuous mapping of X into \([0,1]\) , equal to 1 on \(\operatorname{Supp}(f)\) and with compact support, then \(|f(x)| \leqslant ah(x)\) on \(\operatorname{Supp}(\mu)\) , therefore

\[
| \mu | (| f |) \leqslant a | \mu | (h) \leqslant a \| \mu \|
\]

by Corollary 2; the conclusion then follows from formula (13) of §1, No.~6.

PROPOSITION 9. --- Let \(\mu\) be a positive measure on X; if \(f\) is a function in \(\mathcal{K}_{+}(X)\) such that \(\mu(f)=0\) , then \(f\) is zero on \(\operatorname{Supp}(\mu)\) .

Let \(x\) be a point of \(X\) such that \(f(x) > 0\) ; let us show that \(x\) does not belong to \(\operatorname{Supp}(\mu)\) . For, there exist a compact neighborhood \(V\) of \(x\) and a number \(a > 0\) such that \(f(y) \geqslant a\) on \(V\) . If \(g\) is any continuous function \(\geqslant 0\) with support contained in \(V\) , let us show that \(\mu(g) = 0\) ; indeed, if one sets \(b = \|g\|\) then \(g \leqslant bf/a\) , whence \(\mu(g) \leqslant b\mu(f)/a = 0\) .

PROPOSITION 10. --- Let \(\mu\) be a measure on a locally compact space \(X\) ; for every function \(g \in \mathcal{C}(X; \mathbf{C})\) , the support of the measure \(g \cdot \mu\) is the closure \(T\) of the set of points \(x \in \operatorname{Supp}(\mu)\) such that \(g(x) \neq 0\) .

Set \(S = \operatorname{Supp}(\mu)\) ; let \(x_{0}\) be a point not belonging to T; there exists an open neighborhood V of \(x_{0}\) such that at every point of \(V \cap S\) , g is zero; if \(f \in \mathcal{K}(X; \mathbf{C})\) has support contained in V, then fg is zero on S, therefore (Prop. 8) \(\mu(gf) = 0\) ; in other words, the restriction of \(g \cdot \mu\) to V is zero.

Conversely, assuming that the restriction of \(g \cdot \mu\) to an open neighborhood W of a point \(x_{0} \in X\) is zero, let us show that there does not exist a point of \(W \cap S\) at which g is \(\neq 0\) . Indeed, if there were such a point y, there would exist a compact neighborhood U of y, contained in W, at every point x of which \(g(x) \neq 0\) ; but then every function \(f \in \mathcal{K}(X; \mathbf{C})\) with support contained in U could be written f = gh, where \(h \in \mathcal{K}(X; \mathbf{C})\) has support contained in \(U \subset W\) ; it would then follow that \(\mu(f) = \mu(gh) = 0\) , contrary to the hypothesis \(y \in S\) .

\pandocbounded{\includegraphics[keepaspectratio]{images/part001_a1e78629.jpg}}

Note that T is contained in the intersection of the support S of \(\mu\) and the support of g, but it is not necessarily equal to this intersection. For example, if X = R, \(\mu\) is the Dirac measure at the point 0, and \(g(x) = x\) , then \(g \cdot \mu = 0\) even though the intersection of the supports of g and \(\mu\) reduces to the point 0, thus is nonempty.

COROLLARY. --- In order that the measure \(g \cdot \mu\) be zero, it is necessary and sufficient that g be zero on the support of \(\mu\) .

PROPOSITION 11. --- Every measure with compact support is bounded.

For, \(|\mu|\) is also a measure with compact support, thus we can restrict attention to the case that \(\mu \geqslant 0\) ; if \(h\) is a continuous mapping of \(X\) into {[}0,1{]}, with compact support and equal to 1 on \(\operatorname{Supp}(\mu)\) , then for every function \(f \in \mathcal{K}(X; \mathbf{C})\) one has \(|f(x)| \leqslant \|f\|h(x)\) on \(\operatorname{Supp}(\mu)\) , therefore (Cor. 2 of Prop. 8) \(\mu(|f|) \leqslant \mu(h)\|f\|\) , which proves the proposition (§1, No.~8).

